\documentclass{article}
\setcounter{tocdepth}{4} %Required to insert paragraph in TOC
\setcounter{secnumdepth}{4} %Required to insert paragraph in TOC

\usepackage{graphicx} % Required for inserting images
\graphicspath{ {./img/} }

\usepackage{blindtext}
\usepackage{titlesec}
\usepackage[dvipsnames]{xcolor}
\usepackage{geometry}
\geometry{
	a4paper,
	total={170mm,257mm},
	left=25mm,
	top=20mm,
}

\usepackage{adjustbox}
\usepackage{amsmath}
\usepackage{amsthm}
\usepackage{amssymb}
\usepackage{cancel}
\usepackage{caption}
\usepackage{subcaption}
\usepackage{titlesec}
\usepackage{float}
\usepackage{enumitem}
\usepackage{hyperref}
\usepackage{makecell}
\usepackage{chronology}
\usepackage{cleveref}
\usepackage{fancyhdr}
\usepackage{mathabx}
\usepackage{mathrsfs}
\usepackage{subfiles} %best loaded last in preamble

\pagestyle{fancy}

\setcounter{secnumdepth}{4}

\newcommand{\foo}{\hspace{-2.3pt}$\bullet$ \hspace{5pt}}

\titleformat{\paragraph}
{\normalfont\normalsize\bfseries}{\theparagraph}{1em}{}
\titlespacing*{\paragraph}
{0pt}{3.25ex plus 1ex minus .2ex}{1.5ex plus .2ex}

\newtheorem{definition}{Definizione}[section]
\newtheorem{theorem}{Teorema}[section]
\newtheorem{exmp}{Esempio}[section]

\numberwithin{equation}{subsection}

\title{Appunti Fisica 2}
\author{Mattia Robuschi Caprara}
\date{}

\definecolor{CoverGreen}{RGB}{91, 191, 240}

\begin{document}

	\begin{titlepage}
		
		\pagecolor{CoverGreen}
		
		\vspace{25 mm}
		\begin{center}
			\large
			{\color{black}\textbf{Mattia Robuschi Caprara}} 
		\end{center}
		
		\begin{center}
			\huge
			{\color{black}\textbf{Appunti Fisica 2}}
		\end{center}
		
		\vspace{45 mm}
		
		\begin{figure}[h]
			\centering
			\adjustbox{cfbox=white 5pt 0cm}{\includegraphics[scale=0.3]{0.cover.jpg}}
			%\caption{Caption}
			\label{fig:cover}
		\end{figure}
		
		\thispagestyle{empty} 
	\end{titlepage}

	\newpage
	\pagecolor{white}
	\section*{Prefazione}
	Un riassunto di Fisica 2 (Professor. Massimo Ghidini)
	\tableofcontents
	\newpage
	
	\begin{center}
		\begin{figure}[h]
			\centering
			\includegraphics[scale=0.55]{example-image-a}
			\caption{}
			\label{fig:im-0}
		\end{figure}
	\end{center}
	
	\section{Elettrostatica}
	\subfile{sections/elettrostat}
	
	\section{Elettrodinamica}
	\subfile{sections/elettrodin}
	
	\section{Magnetostatica}
	\subfile{sections/magnetostat}
	
	\section{Induzione Elettromagnetica}
	\subfile{sections/induzElettromag}
	
	\section{Onde Elettromagnetiche}
	\subfile{sections/onde}
\end{document}